\begin{figure}[t]
\centering
% Shared color TikZ styles for the paper figures.
\colorlet{paperBlue}{blue!9!white}
\colorlet{paperBlueLine}{blue!55!black}
\colorlet{paperTeal}{teal!10!white}
\colorlet{paperTealLine}{teal!55!black}
\colorlet{paperAmber}{orange!15!white}
\colorlet{paperAmberLine}{orange!60!black}
\colorlet{paperRose}{red!9!white}
\colorlet{paperRoseLine}{red!55!black}
\colorlet{paperViolet}{violet!10!white}
\colorlet{paperVioletLine}{violet!55!black}
\colorlet{paperInk}{black!72}

\tikzset{
  paperfig/.style={
    font=\small,
    >=Stealth,
    line cap=round,
    line join=round
  },
  paperline/.style={
    -{Stealth[length=2.1mm,width=1.6mm]},
    line width=0.58pt,
    draw=paperInk
  },
  paperdash/.style={
    paperline,
    dashed,
    draw=black!48
  },
  paperbox/.style={
    draw=black!58,
    line width=0.58pt,
    rounded corners=2pt,
    fill=white,
    align=center,
    inner xsep=5pt,
    inner ysep=5pt,
    minimum height=8mm
  },
  paperdata/.style={
    paperbox,
    fill=paperBlue,
    draw=paperBlueLine
  },
  paperproc/.style={
    paperbox,
    fill=paperTeal,
    draw=paperTealLine
  },
  papermodel/.style={
    paperbox,
    fill=paperAmber,
    draw=paperAmberLine,
    line width=0.65pt
  },
  paperout/.style={
    paperbox,
    fill=paperRose,
    draw=paperRoseLine,
    line width=0.65pt,
    font=\small\bfseries
  },
  paperstate/.style={
    paperbox,
    fill=paperViolet,
    draw=paperVioletLine
  },
  papernote/.style={
    font=\scriptsize,
    align=center,
    text=black!70,
    fill=white,
    inner sep=1.4pt
  },
  paperaxis/.style={
    line width=0.45pt,
    draw=black!72
  },
  papergrid/.style={
    line width=0.25pt,
    draw=black!15
  },
  paperplot/.style={
    line width=0.95pt,
    draw=paperBlueLine
  },
  paperplotA/.style={
    line width=0.95pt,
    draw=paperBlueLine
  },
  paperplotB/.style={
    line width=0.95pt,
    draw=paperAmberLine
  },
  papermark/.style={
    circle,
    draw=paperBlueLine,
    fill=white,
    line width=0.55pt,
    inner sep=1.6pt
  },
  papermarkA/.style={
    papermark,
    draw=paperBlueLine,
    fill=paperBlue
  },
  papermarkB/.style={
    papermark,
    draw=paperAmberLine,
    fill=paperAmber
  }
}

\resizebox{0.98\linewidth}{!}{%
\begin{tikzpicture}[paperfig, node distance=5mm and 7mm]
  \node[paperdata, text width=31mm] (task) {任务输入\\\scriptsize 题面、签名、测试};
  \node[papermodel, below=of task, text width=31mm] (single) {单步\\\scriptsize 一次生成};
  \node[paperproc, below=of single, text width=31mm] (repair) {单步+修补\\\scriptsize 执行反馈修补};
  \node[paperproc, below=of repair, text width=31mm] (route) {分流+重链与救援\\\scriptsize 结构路由与失败闭环};
  \node[paperproc, below=of route, text width=31mm] (cascade) {级联\\\scriptsize 候选扩展与筛选};
  \node[paperout, below=of cascade, text width=31mm] (eval) {统一评测\\\scriptsize pass@1 / calls / tokens};

  \draw[paperline] (task) -- (single);
  \draw[paperline] (single) -- (repair);
  \draw[paperline] (repair) -- (route);
  \draw[paperline] (route) -- (cascade);
  \draw[paperline] (cascade) -- (eval);

  \node[papernote, right=8mm of repair] (mid) {中预算\\执行反馈};
  \node[papernote, right=8mm of route] (heavy) {结构预算\\复杂样本};
  \node[papernote, right=8mm of cascade] (high) {高预算\\候选覆盖};
  \draw[paperdash] (mid) -- (repair);
  \draw[paperdash] (heavy) -- (route);
  \draw[paperdash] (high) -- (cascade);

  \draw[paperdash] (route.west) to[bend left=28]
    node[papernote, left, text width=18mm] {失败信息\\压缩回注}
    (repair.west);
\end{tikzpicture}%
}
\caption{预算分档编排整体框架}
\label{fig:overall_framework}
\end{figure}
